\subsubsection{Đối chiếu với Project Baseline V1.0}
\label{sec:t3-actual}

\subsubsubsection{Phương pháp xác định số liệu thực tế}
\label{sec:t3-effort-method}

\textbf{Ngày thực tế} được lấy từ bằng chứng kiểm chứng được: commit đầu tiên chứa đầy đủ chức năng, ngày merge, hoặc ngày của biên bản kiểm thử trong repository.

\textbf{Giờ công thực tế} được tổng hợp từ \textbf{bảng chấm công} của nhóm: mỗi thành viên ghi số giờ đã làm cho từng hoạt động của baseline và cho từng tuần, kèm nguồn đối chiếu như commit, tin nhắn nhóm hoặc lịch học. PM tổng hợp, kiểm tra tổng theo hoạt động khớp với tổng theo tuần của từng người, và đối chiếu các hoạt động lập trình với lịch sử commit. Các việc không có trong baseline được ghi thành dòng ``phát sinh'' riêng.

\subsubsubsection{Ngày thực tế và công sức của các hoạt động chính}

Độ lệch tính bằng ngày lịch, dương là trễ, âm là sớm.

\begingroup\footnotesize
\renewcommand{\arraystretch}{1.15}
\begin{longtable}{@{}p{.07\linewidth}p{.11\linewidth}p{.20\linewidth}p{.06\linewidth}p{.06\linewidth}p{.06\linewidth}p{.34\linewidth}@{}}
\caption{Kế hoạch V1.0 so với thực tế (đến 27/11/2026)}\label{tab:t3-actual}\\
\toprule
\textbf{Mã} & \textbf{Lịch KH} & \textbf{Hoàn thành TT (bằng chứng)} & \textbf{Lệch ngày} & \textbf{Giờ KH} & \textbf{Giờ TT} & \textbf{Giải thích}\\\midrule\endfirsthead
\toprule
\textbf{Mã} & \textbf{Lịch KH} & \textbf{Hoàn thành TT (bằng chứng)} & \textbf{Lệch ngày} & \textbf{Giờ KH} & \textbf{Giờ TT} & \textbf{Giải thích}\\\midrule\endhead
\midrule\multicolumn{7}{r}{\textit{Còn tiếp}}\\\endfoot
\bottomrule\endlastfoot
A03 & 29--30/10 & 03/11 (\texttt{b57b1aa}, smoke test) & +4 & 8,0 & 10,0 & Thêm thời gian xử lý sự cố Docker, socket Java, múi giờ JDBC (AR3).\\
A05.1 & 02--04/11 & 03/11 (\texttt{2ed54e3}) & $-$1 & 18,0 & 14,0 & Ít việc hơn dự kiến do mỗi tài khoản chỉ có một vai trò cố định (CR-01). Merge vào nhánh chính 12/11 do chờ review.\\
A05.2 & 30/10--03/11 & 06/11 (\texttt{aa2494d}) & +3 & 18,0 & 31,0 & Giao diện làm lại theo nhận diện HCMUT, chỉnh sửa và refactor nhiều vòng.\\
A06 & 05--09/11 & 07/11 (regression) & $-$2 & 10,0 & 10,0 & Đúng kế hoạch; 208/209 ca đạt.\\
A07.1 & 09--11/11 & 12/11 (\texttt{43f8c4c}) & +1 & 22,0 & 15,0 & Dùng lại kiến trúc phân tầng; làm cùng A07.2, A07.3.\\
A07.2 & 12--16/11 & 12/11 (\texttt{43f8c4c}) & $-$4 & 14,0 & 9,0 & Logic đôn chờ dùng chung transaction với A07.1.\\
A07.3 & 17/11 & 12/11 (\texttt{43f8c4c}) & $-$5 & 6,0 & 5,5 & Outbox làm cùng đăng ký để rollback đồng bộ.\\
A08 & 18--19/11 & 18/11 (PR \#3) & $-$1 & 13,0 & 12,0 & Frontend I2 và kiểm thử I2.\\
A09.1 & 21--24/11 & 23/11 (\texttt{26efcbb}) & $-$1 & 12,0 & 16,0 & Thêm việc sau rà soát bảo mật: tách xem trước và xác nhận, kiểm tra ở database (CR-07).\\
A09.2 & 20--23/11 & 19/11 (\texttt{e8f23a2}) & $-$4 & 10,0 & 10,0 & Đúng kế hoạch; dùng chung API với A09.1.\\
A10.1 & 20--21/11 & 22/11 (\texttt{98f6024}) & +1 & 13,0 & 13,0 & Đúng kế hoạch.\\
A10.2 & 23/11 & 22/11 (\texttt{98f6024}) & $-$1 & 5,0 & 5,0 & Chỉ CSV (CR-02); phần tiết kiệm được dùng cho chống CSV injection, giới hạn 10.000 dòng.\\
A11 & 24--25/11 & 25/11 (\texttt{727f379}, deploy) & 0 & 7,0 & 8,0 & Thêm thời gian sửa build (Maven 429) và chuyển tiếp API trên Vercel (CR-04).\\
A12.1 & 26--27/11 & Thu gọn (CR-05) & -- & 14,0 & 3,5 & Chỉ kiểm tra chi tiết điểm danh (\texttt{bd77472}) và kiểm tra cloud 27/11; không chạy end-to-end trên cloud.\\
RA01 & 09/11 & 12/11 (I2\_RESULTS) & +3 & 2,5 & 2,5 & Thử 6 sinh viên tranh 1 chỗ trong bộ test I2.\\
RA02 & 05/11 & Không thực hiện & -- & 2,5 & 0,0 & Thay bằng triển khai Render/Vercel/Neon (CR-04).\\
RA05 & 23/11 & 18/11 (I3\_RESULTS) & $-$5 & 1,9 & 1,0 & Đo trên 2.000 đăng ký khi audit I3: báo cáo 0,066 s, CSV 0,909 s.\\
RA06 & 09/11 & 12/11 & +3 & 1,2 & 1,2 & Đối chiếu yêu cầu I2 với Client Interaction Log khi bắt đầu I2.\\
\end{longtable}
\endgroup

\subsubsubsection{Biểu đồ theo dõi tiến độ (tracking Gantt)}

Thanh xám là lịch kế hoạch V1.0; hình thoi là ngày hoàn thành thực tế có bằng chứng (Bảng \ref{tab:t3-actual}). Hình thoi nằm bên phải thanh là trễ, bên trái là sớm. Đường đỏ đánh dấu ngày báo cáo 27/11.

\begin{figure}[H]
\centering
\begin{adjustbox}{max width=\linewidth}
\begin{tikzpicture}[x=0.42cm, y=0.5cm]
  \draw[gray!40] (0,0.5) -- (0,-14.2);
  \node[font=\tiny, anchor=south] at (3.5,0.5) {T1};
  \draw[gray!40] (7,0.5) -- (7,-14.2);
  \node[font=\tiny, anchor=south] at (10.5,0.5) {T2};
  \draw[gray!40] (14,0.5) -- (14,-14.2);
  \node[font=\tiny, anchor=south] at (17.5,0.5) {T3};
  \draw[gray!40] (21,0.5) -- (21,-14.2);
  \node[font=\tiny, anchor=south] at (24.5,0.5) {T4};
  \draw[gray!40] (28,0.5) -- (28,-14.2);
  \node[font=\tiny, anchor=south] at (31.5,0.5) {T5};
  \draw[gray!40] (35,0.5) -- (35,-14.2);
  \node[font=\tiny, anchor=south] at (38.5,0.5) {T6};
  \node[font=\tiny, anchor=north, gray] at (0,-14.2) {26/10};
  \node[font=\tiny, anchor=north, gray] at (7,-14.2) {02/11};
  \node[font=\tiny, anchor=north, gray] at (14,-14.2) {09/11};
  \node[font=\tiny, anchor=north, gray] at (21,-14.2) {16/11};
  \node[font=\tiny, anchor=north, gray] at (28,-14.2) {23/11};
  \node[font=\tiny, anchor=north, gray] at (35,-14.2) {30/11};
  \node[font=\scriptsize, anchor=east] at (-0.3,-1) {A03};
  \fill[gray!45] (3,-1.28) rectangle (5,-0.72);
  \node[diamond, fill=blue!70!black, inner sep=1.6pt] at (8.5,-1) {};
  \node[font=\scriptsize, anchor=east] at (-0.3,-2) {A05.1};
  \fill[gray!45] (7,-2.28) rectangle (10,-1.72);
  \node[diamond, fill=blue!70!black, inner sep=1.6pt] at (8.5,-2) {};
  \node[font=\scriptsize, anchor=east] at (-0.3,-3) {A05.2};
  \fill[gray!45] (4,-3.28) rectangle (9,-2.72);
  \node[diamond, fill=blue!70!black, inner sep=1.6pt] at (11.5,-3) {};
  \node[font=\scriptsize, anchor=east] at (-0.3,-4) {A06};
  \fill[gray!45] (10,-4.28) rectangle (15,-3.72);
  \node[diamond, fill=blue!70!black, inner sep=1.6pt] at (12.5,-4) {};
  \node[font=\scriptsize, anchor=east] at (-0.3,-5) {A07.1};
  \fill[gray!45] (14,-5.28) rectangle (17,-4.72);
  \node[diamond, fill=blue!70!black, inner sep=1.6pt] at (17.5,-5) {};
  \node[font=\scriptsize, anchor=east] at (-0.3,-6) {A07.2};
  \fill[gray!45] (17,-6.28) rectangle (22,-5.72);
  \node[diamond, fill=blue!70!black, inner sep=1.6pt] at (17.5,-6) {};
  \node[font=\scriptsize, anchor=east] at (-0.3,-7) {A07.3};
  \fill[gray!45] (22,-7.28) rectangle (23,-6.72);
  \node[diamond, fill=blue!70!black, inner sep=1.6pt] at (17.5,-7) {};
  \node[font=\scriptsize, anchor=east] at (-0.3,-8) {A08};
  \fill[gray!45] (23,-8.28) rectangle (25,-7.72);
  \node[diamond, fill=blue!70!black, inner sep=1.6pt] at (23.5,-8) {};
  \node[font=\scriptsize, anchor=east] at (-0.3,-9) {A09.1};
  \fill[gray!45] (26,-9.28) rectangle (30,-8.72);
  \node[diamond, fill=blue!70!black, inner sep=1.6pt] at (28.5,-9) {};
  \node[font=\scriptsize, anchor=east] at (-0.3,-10) {A09.2};
  \fill[gray!45] (25,-10.28) rectangle (29,-9.72);
  \node[diamond, fill=blue!70!black, inner sep=1.6pt] at (24.5,-10) {};
  \node[font=\scriptsize, anchor=east] at (-0.3,-11) {A10.1};
  \fill[gray!45] (25,-11.28) rectangle (27,-10.72);
  \node[diamond, fill=blue!70!black, inner sep=1.6pt] at (27.5,-11) {};
  \node[font=\scriptsize, anchor=east] at (-0.3,-12) {A10.2};
  \fill[gray!45] (28,-12.28) rectangle (29,-11.72);
  \node[diamond, fill=blue!70!black, inner sep=1.6pt] at (27.5,-12) {};
  \node[font=\scriptsize, anchor=east] at (-0.3,-13) {A11};
  \fill[gray!45] (29,-13.28) rectangle (31,-12.72);
  \node[diamond, fill=blue!70!black, inner sep=1.6pt] at (30.5,-13) {};
  \node[font=\scriptsize, anchor=east] at (-0.3,-14) {A12.1};
  \fill[gray!45] (31,-14.28) rectangle (33,-13.72);
  \draw[red, thick] (32.5,0.5) -- (32.5,-14.2);
  \node[font=\tiny, red, anchor=south west] at (32.6,0.5) {27/11};
\end{tikzpicture}
\end{adjustbox}
\vspace{4pt}
\caption{Lịch kế hoạch V1.0 và ngày hoàn thành thực tế các hoạt động phát triển}\label{fig:t3-tracking-gantt}
\end{figure}

Biểu đồ cho thấy rõ hai giai đoạn: đến CP1, các hoạt động trễ do A03 (hình thoi bên phải thanh); từ I2, hầu hết hoạt động hoàn thành sớm hơn kế hoạch, riêng A07.1 và A10.1 trễ 1 ngày. Các hoạt động A07.1--A07.3 có cùng một ngày hoàn thành vì được gộp làm chung (Mục \ref{sec:t3-actual}).


\subsubsubsection{Công việc phát sinh ngoài baseline}

\begin{table}[H]
\centering\small
\caption{Công việc không có trong Project Baseline V1.0}\label{tab:t3-unplanned}
\begin{tabularx}{\linewidth}{@{}p{.10\linewidth}p{.30\linewidth}X p{.09\linewidth}@{}}
\toprule
\textbf{Ngày} & \textbf{Công việc} & \textbf{Lý do thực hiện} & \textbf{Giờ TT}\\\midrule
12/11 & Tái cấu trúc backend phân tầng (\texttt{4d56706}) & Tránh viết lại tầng dữ liệu khi thêm I2/I3. Phần lớn là di chuyển mã nên chỉ tính một nửa số dòng thay đổi. & 12,0\\
25/11 & Rà soát hệ thống: giới hạn đăng nhập sai, quản lý phiên bản database, sao lưu và khôi phục, kiểm thử tự động (\texttt{275bfb9}) & Điều kiện cần để đưa lên cloud an toàn; không có hoạt động tương ứng trong A11. & 10,0\\
26/11 & Chỉnh mẫu email (\texttt{13c81af}) & Cải thiện nội dung thư sau khi xem thư thật. & 2,0\\
\midrule
\multicolumn{3}{@{}r}{\textbf{Tổng công việc phát sinh}} & \textbf{24,0}\\
\bottomrule
\end{tabularx}
\end{table}

\subsubsubsection{Giải thích các sai lệch chính}

\begin{enumerate}
    \item \textbf{Giai đoạn đầu trễ và tốn công hơn.} A03 trễ 4 ngày kéo A05.2 trễ 3 ngày. A05.2 vượt 13 giờ (+72\%), là sai lệch lớn nhất, vì ước lượng chỉ tính một lượt giao diện, không tính phần thiết kế lại theo nhận diện HCMUT.
    \item \textbf{I2 hoàn thành sớm và ít công hơn.} A07.1--A07.3 dùng 29,5 giờ so với 42 giờ kế hoạch ($-$30\%). Nguyên nhân chính là quyết định fast-tracking sau CR-05 (Mục \ref{sec:t3-change}); ngoài ra kiến trúc phân tầng vừa tái cấu trúc giúp thêm tính năng nhanh; ba hoạt động được làm chung vì cùng một transaction. Nhóm có dùng công cụ AI hỗ trợ lập trình, được khai báo ở Mục \ref{sec:ai-declaration}. Tuy nhiên phần tiết kiệm này gần như bị bù lại bởi 12 giờ tái cấu trúc ngoài kế hoạch.
    \item \textbf{I3 sát ước lượng.} A09--A10 dùng 44 giờ so với 40 giờ (+10\%), phần vượt đến từ audit bảo mật QR.
    \item \textbf{Tích hợp dồn về cuối.} Mỗi increment được merge vào nhánh kế tiếp thay vì nhánh chính, nên đến 25/11 mới tích hợp toàn bộ, phát sinh 10 giờ rà soát ngoài kế hoạch.
    \item \textbf{Tổng thể.} Đến 27/11, công sức thực tế là 306,2 giờ so với 302,2 giờ kế hoạch cho T1--T5 (+1,3\%). Công việc phát sinh (24 giờ) được bù bởi phần tiết kiệm ở I2 và các hoạt động chưa bắt đầu (A12.2, A14.1).
\end{enumerate}

\subsubsubsection{Chỉ số giá trị thu được (EVM) tại ngày 27/11/2026}

Giá trị tính bằng giờ công kế hoạch (đơn giá 78.000 VNĐ/giờ):
\begin{itemize}
    \item $PV$ = 293,6 giờ (22.900.800 VNĐ): tổng giờ kế hoạch của các hoạt động lẽ ra đã xong trước 27/11, cộng phần quản lý G05/G06 theo thời gian.
    \item $EV$ = 280,6 giờ (21.886.800 VNĐ): tổng giờ kế hoạch của các hoạt động đã xong; A12.1 được tính 25\% vì mới thực hiện một phần; RA02 không tính vì không thực hiện.
    \item $AC$ = 306,2 giờ (23.883.600 VNĐ): công sức thực tế theo bảng chấm công.
\end{itemize}

\begin{center}
$SPI = EV/PV = 280{,}6 / 293{,}6 = 0{,}96$ \qquad $CPI = EV/AC = 280{,}6 / 306{,}2 = 0{,}92$
\end{center}


\begin{figure}[H]
\centering
\begin{tikzpicture}
\begin{axis}[xmin=0.7, xmax=5.3, ymin=0, ymax=340,
    xtick={1,2,3,4,5}, xticklabels={T1,T2,T3,T4,T5},
    ylabel={Giờ công tích lũy},
    width=.92\linewidth, height=6.2cm,
    yticklabel style={/pgf/number format/use comma},
    every node near coord/.append style={font=\tiny, /pgf/number format/use comma, /pgf/number format/fixed},
    legend style={font=\scriptsize, draw=none, at={(0.5,1.03)}, anchor=south, legend columns=-1},
    tick label style={font=\scriptsize}, label style={font=\small},
    ymajorgrids, grid style={gray!20}]
\addplot[thick, gray!70, mark=square*, mark options={fill=gray!45}] coordinates {(1,63.5) (2,127.5) (3,188.6) (4,248.8) (5,302.2)};
\addplot[thick, blue!75, mark=*] coordinates {(1,56.5) (2,136.0) (3,205.7) (4,270.7) (5,306.2)};
\addplot[only marks, mark=diamond*, mark size=4pt, green!50!black] coordinates {(5,280.6)};
\addplot[dashed, gray!60, domain=0.7:5.3, samples=2] {331.6};
\node[font=\tiny, anchor=south west, gray] at (axis cs:0.75,331.6) {BAC = 331,6};
\legend{Kế hoạch (PV), Thực tế (AC), Giá trị đạt được (EV) tại 27/11}
\end{axis}
\end{tikzpicture}
\vspace{6pt}
\caption{Đường cong S: công sức tích lũy kế hoạch và thực tế}\label{fig:t3-scurve}
\end{figure}

Đường thực tế nằm trên đường kế hoạch từ T2 đến T4 (tốn công sớm hơn dự kiến) và gặp lại ở T5. Điểm EV nằm dưới cả hai đường, phản ánh $SPI$ và $CPI$ nhỏ hơn 1 tại ngày đo.

$SPI = 0{,}96$ cho thấy dự án chậm nhẹ tại ngày đo, chủ yếu do A12 (integration test) mới thực hiện một phần; sau đó A12 được thu gọn theo CR-05 thay vì hoàn thành đầy đủ. Các chức năng phần mềm đều đã hoàn thành. $CPI = 0{,}92$ cho thấy mỗi giờ công tạo ra khoảng 0,92 giờ giá trị kế hoạch, do 24 giờ công việc phát sinh. Nếu hiệu suất này giữ nguyên, ước tính tổng công sức khi hoàn thành là $EAC = BAC / CPI = 331{,}6 / 0{,}92 \approx 360$ giờ, vẫn nằm trong quỹ dự phòng 63,4 giờ (tổng năng lực 395 giờ).
